\section{Scope, provenance, and the exact questions continued}
\label{sec:scope}

The starting point is the current collective manuscript
\emph{Polylogarithms and their Arithmetic Bridges} and the six
archives in the repository's incoming directory at the title-page
commit \cite{Repo}. The source was obtained as a pinned Git checkout;
the incoming archives were extracted into separate directories outside
that checkout. The research package is an addition for later
integration. It does not change the canonical manuscript or the
incoming archive contents.

The two newest digamma reports are particularly relevant.
\emph{Canonical Digamma and Fifth Moment} proves a finite, but
generally longer, initial-moment reconstruction theorem and conjectures
that the first $q$ moments already suffice for unequal even orders
\cite[Sections 7 and 9, \texttt{conj:firstq}]{CanonicalFifth}.
\emph{Higher Digamma Powers and Ordered Gamma Sums} asks for an
organized sixth-power reduction and for a differential-polynomial
integration calculus \cite[Further exact research questions]{HigherPowers}.
The previous coincident/directional-zeta report supplies the complete
quadratic ray formula for Tornheim functions, while leaving cubic
arithmetic data unresolved \cite[Section 6]{DirectionalPrior}.

\subsection{Main contributions and their limits}

The first contribution resolves an infinite subfamily of a stated
conjecture, rather than adding another finite rank experiment.
For the pair $(2,2b)$, an explicit Pad\'e calculation evaluates the
coefficient determinant and a rational extraction row. Classical
Bernoulli Hankel machinery is reproved as needed; the mixed matrix
perturbation and its application to the moment conjecture are the
new continuation \cite{DilcherJiu2020}. The resulting sums and the
polylogarithmic boundary family are identities at every $b\ge2$.
A rank-two version proves the entire $(4,2b)$ family. A congruence
modulo four and a finite polynomial certificate prove $(6,2b)$ as well. For every
even pair $(2a,2b)$ the remaining determinant is reduced explicitly to
size $a$; the cases $8\le p<q$ remain unresolved here.

The second contribution is an exact cancellation theorem for third
spectral derivatives. The spectral origin is not a regular point of
the three-variable Tornheim function. We first isolate its polar
part, then differentiate specified holomorphic ray restrictions.
The elementary axis and quasi-shuffle restrictions determine enough
of the remaining Taylor coefficients to yield a complete two-coordinate
formula and explicit combinations containing ordinary zeta derivatives
alone. The formal ambiguity theorem states precisely what these
particular local restrictions leave undetermined. It is not a
transcendence or arithmetic-independence theorem.

The third contribution answers the algebraic part of the proposed
differential-polynomial classification. Its partition basis gives
an exact finite algorithm for reduction modulo total derivatives.
This is a constructive specialization of classical differential
algebra \cite{OlverShakiban1978}, with a new application to the current
Gamma-normalized finite-part problem. The normalizations at degrees
six and seven are computed explicitly. A remaining mixed integral is displayed,
not silently declared independent or replaced by an unevaluated
quantity under a new name.

The principal status changes are summarized in
Table~\ref{tab:status}. All references to novelty in this article
mean an addition relative to the inspected ProveIt material.
A complete historical priority claim would require a broader
comparison with all specializations of the classical literature.

\begin{table}[htbp]
\centering\small
\caption{Source questions and the result established here.}
\label{tab:status}
\begin{tabular}{@{}p{.29\textwidth}p{.42\textwidth}p{.22\textwidth}@{}}
\toprule
Question & Proved result & Remaining scope\\
\midrule
Optimal first-$q$ prefix & Complete proofs for smaller order $2$, $4$, or $6$; a fixed-size determinant for every even pair & Unequal $8\le p<q$\\[5pt]
Higher Tornheim ray derivatives & All cubic directions; exact cancellation criterion and an all-degree local kernel & Arithmetic values of the two cubic coordinates\\[5pt]
Differential digamma classification & Canonical partition basis, bidegree dimensions, Gamma parity and endpoints & The image after additional spectral relations\\[5pt]
Sixth and seventh moments & Canonical sixth formula with one mixed moment; explicit seventh formula & Spectral evaluation of the mixed integrals\\[5pt]
Short Gaussian evaluations & No status change & $S_6$ and revised $S_8$\\
\bottomrule
\end{tabular}
\end{table}

\subsection{Conventions that affect the constants}

For $r\ge1$, set $H_n^{(r)}=\sum_{k=1}^n k^{-r}$ and
$H_0^{(r)}=0$. Bernoulli numbers and polynomials use
\[
 \frac{te^{xt}}{e^t-1}=\sum_{n\ge0}B_n(x)\frac{t^n}{n!},
 \qquad B_n=B_n(0),\qquad B_1=-\frac12.
\]
The digamma function is $\psi(x)=\Gamma'(x)/\Gamma(x)$.
The generalized Stieltjes convention is
\begin{equation}\label{eq:stieltjes-convention}
 \zeta(1+u,a)=\frac1u+
       \sum_{j\ge0}\frac{(-1)^j\gamma_j(a)}{j!}u^j,
 \qquad \re a>0.
\end{equation}
We abbreviate $\gamma=\gamma_0(1)$ and
$\gamma_j=\gamma_j(1)$. Derivatives $\zeta^{(j)}(s)$ are always
with respect to the spectral variable, not to a Hurwitz shift.
Standard Gamma, polygamma, and Hurwitz formulas are used with these
normalizations \cite{DLMF}.

If $f$ is meromorphic near $0$ and regular on $(0,1]$, then
\begin{equation}\label{eq:fp-convention}
 \FP\int_0^1 f(x)\dd x
 :=\CT_{\eps\downarrow0}\int_{\eps}^1 f(x)\dd x.
\end{equation}
The constant term discards negative powers of $\eps$ and logarithms
but retains the constant in the original coordinate $x$.
In particular, for a meromorphic primitive $g$,
\begin{equation}\label{eq:endpoint-law-intro}
 \FP\int_0^1 g'(x)\dd x=g(1)-\CT_{x=0}g(x).
\end{equation}
This convention is distinct from the centered asymptotic moment
continuation in Section~\ref{sec:tail-foundation}; the latter is
defined explicitly before it is used. Both conventions eventually
lead to ordinary convergent identities after their stated subtractions.
